\chapter*{ABSTRACT}

Emergency departments in low- and middle-income countries (LMICs) face overcrowding, delayed first assessment, and a failure we call \emph{acuity discontinuity}: the South African Triage Scale (SATS) colour assigned at the desk often stops guiding care after the patient leaves triage. This thesis presents a hospital chatbot that maps an unstructured chief complaint to a SATS acuity level and colour, then attaches a concrete clinical pathway---destination, target response time, and actions.

\textbf{Part~1} implements a Digital Acuity Engine: a medical gate filters non-clinical text; optional OpenMed named-entity enrichment surfaces diseases and drugs; a fine-tuned BioBERT classifier predicts acuity levels \(1\)--\(5\); a fixed map \(f_{\mathrm{SATS}}\) yields Red, Orange, Yellow, or Green. On an \(8{,}000\)-sample stratified holdout, baseline BioBERT achieved \(99.92\%\) colour-routing accuracy and \(98.14\%\) Level-1 (Red) recall.

\textbf{Part~2} defines protocol-mapped pathways \(P(C)\) for each colour (including Green diversion from the acute ED and a Blue dignity protocol). The architecture deliberately omits Bayesian fusion with Triage Early Warning Score (TEWS) vitals at chatbot intake: BioBERT already performs end-to-end text-to-acuity inference; pathways attach deterministically to colour. TEWS remains nurse-side or future work.

The contribution is a defensible LMIC chatbot design and prototype (Curatio) that answers whether NLP can drive colour-coded clinical pathways without requiring a full TEWS--Bayesian stack before the triage desk.

\vspace{0.5cm}
\noindent\textbf{Keywords:} SATS; clinical pathways; BioBERT; triage chatbot; LMIC; acuity discontinuity.
